% notes-general-formalization.tex
% Topic: General Formalization

\begin{remark*}[Working posture]
Formalization is not transcription. A textbook definition supplies the
mathematical target; a Lean development must also choose the representation,
the library interface, the available theorem names, and the shape of the proof
state after each tactic.
\end{remark*}

\begin{remark*}[Library-first formalization]
When a concept is already present in mathlib, the first task is to locate and
understand the existing API. Re-declaring the concept locally usually creates
duplicate vocabulary and unnecessary conversion work. The better workflow is:
\begin{enumerate}
  \item identify the existing typeclass, structure, notation, or theorem;
  \item state the mathematical reminder in the notes without extracting it as a
        new canonical definition;
  \item write small examples that use the existing API;
  \item record the proof-state obstacles that explain why each theorem or
        tactic is needed.
\end{enumerate}
\end{remark*}

\begin{remark*}[Reading proof branches]
When a Lean proof splits an if-and-only-if statement, the command
\texttt{constructor} creates two goals in order: first the forward implication
and then the reverse implication. For a goal $A\leftrightarrow B$, the first
branch proves $A\to B$ and the second branch proves $B\to A$.

This matters because a proof can be mathematically sensible and still be in
the wrong branch. Temporary comments such as \texttt{-- =>} and
\texttt{-- <=} are useful labels while the proof is being written.
\end{remark*}

\begin{remark*}[Inspecting the proof state]
During exploration, \texttt{trace\_state} can be placed inside a proof to print
the current local context and goal. It is useful when the available
hypotheses do not match the mental picture of the proof. It should be removed
from finished Lean code.
\end{remark*}
